\section{Euler's pentagonal number theorem and Jacobi's triple product identity}

\subsection{Euler's pentagonal number theorem}

\begin{definition}
\label{def.pars.pent-num}
\lean{AlgebraicCombinatorics.pentagonalNumber}
\leantarget
\leanok
For any $k\in\mathbb{Z}$, define a nonnegative
integer $w_{k}\in\mathbb{N}$ by
\[
w_{k}=\dfrac{\left(  3k-1\right)  k}{2}.
\]
This is called the $k$\emph{-th pentagonal number}.
\end{definition}

\begin{theorem}
[Euler's pentagonal number theorem]
\label{thm.pars.pent}
\lean{AlgebraicCombinatorics.euler_pentagonal_number_theorem}
\leantarget
\leanok
We have
\[
\prod_{k=1}^{\infty}\left(  1-x^{k}\right)  =\sum_{k\in\mathbb{Z}}\left(
-1\right)  ^{k}x^{w_{k}}.
\]
\end{theorem}

\begin{proof}
Set $q=x^{3}$ and $z=-x$ in Theorem \ref{thm.pars.jtp2}. (This means that we apply
Theorem \ref{thm.pars.jtp2} to $a=3$ and $b=1$ and $u=1$ and $v=-1$.) We get
\begin{align}
&  \prod_{n>0}\left(  \left(  1+\left(  x^{3}\right)  ^{2n-1}\left(
-x\right)  \right)  \left(  1+\left(  x^{3}\right)  ^{2n-1}\left(  -x\right)
^{-1}\right)  \left(  1-\left(  x^{3}\right)  ^{2n}\right)  \right)
\nonumber\\
&  =\sum_{\ell\in\mathbb{Z}}\left(  x^{3}\right)  ^{\ell^{2}}\left(
-x\right)  ^{\ell}. \label{pf.thm.pars.pent.1}
\end{align}
The left hand side of this equality simplifies as follows:
\begin{align*}
&  \prod_{n>0}\left(  \left(  1-x^{6n-2}\right)  \left(
1-x^{6n-4}\right)  \left(  1-x^{6n}\right)  \right) \\
&  =\prod_{n>0}\left(  \left(  1-\left(  x^{2}\right)
^{3n-1}\right)  \left(  1-\left(  x^{2}\right)
^{3n-2}\right)  \left(  1-\left(  x^{2}\right)
^{3n}\right)  \right) \\
&  =\prod_{k>0}\left(  1-\left(  x^{2}\right)  ^{k}\right)  ,
\end{align*}
since each positive integer $k$ can be uniquely represented as $3n-1$ or
$3n-2$ or $3n$ for some positive integer $n$.

Comparing this with (\ref{pf.thm.pars.pent.1}), we obtain
\begin{align}
&  \prod_{k>0}\left(  1-\left(  x^{2}\right)  ^{k}\right) \nonumber\\
&  =\sum_{\ell\in\mathbb{Z}}\left(  -1\right)  ^{\ell}x^{3\ell
^{2}+\ell}=\sum_{\ell\in\mathbb{Z}}\left(  -1\right)  ^{\ell}x^{2w_{-\ell}} \nonumber\\
&  =\sum_{k\in\mathbb{Z}}\left(  -1\right)  ^{k}\left(  x^{2}\right)  ^{w_{k}}
\label{pf.thm.pars.pent.2}
\end{align}
(here, we have substituted $k$ for $-\ell$ in the sum, and used the
identity $3\ell^{2}+\ell = 2w_{-\ell}$).

Now, we ``substitute $x$ for $x^{2}$'' in
this equality using Lemma \ref{lem.fps.fxx=gxx}. As a result, we obtain
\[
\prod_{k>0}\left(  1-x^{k}\right)  =\sum_{k\in\mathbb{Z}}\left(  -1\right)
^{k}x^{w_{k}}.
\]
This is Euler's pentagonal number theorem (Theorem \ref{thm.pars.pent}).

The Lean proof proceeds by first proving the identity over $\mathbb{Q}$
(using Lemma~\ref{lem.jtp.euler-pentagonal-rat},
Lemma~\ref{lem.jtp.jacobiLHS-eq-eulerProduct-sq},
Lemma~\ref{lem.jtp.jacobiRHS-eq-pentagonalSeries-sq}, and
Lemma~\ref{lem.fps.fxx=gxx}), and then transferring to $\mathbb{Z}$
by comparing coefficients via the ring homomorphism $\mathbb{Z}\to\mathbb{Q}$
(using Lemma~\ref{lem.jtp.eulerProduct-coeff-map}).
\end{proof}

\subsubsection{Pentagonal number properties}

\begin{lemma}
\label{lem.pent.injective}
\lean{AlgebraicCombinatorics.pentagonalNumber_injective}
\leanhelper
\leanok
The function $k\mapsto w_k$ is injective: if $w_j = w_k$ then $j = k$.
\end{lemma}

\begin{proof}
\leanok
By monotonicity on the positive integers (strictly increasing) and on the negative integers (also strictly increasing in absolute value), and the fact that
$w_k = w_j$ forces $k, j$ to have the same sign.
\end{proof}

\begin{lemma}
\label{lem.pent.below-finite}
\lean{AlgebraicCombinatorics.pentagonal_below_finite}
\leanhelper
\leanok
For any $n\in\mathbb{N}$, the set $\{k\in\mathbb{Z} : w_k < n\}$ is finite.
\end{lemma}

\begin{proof}
\leanok
Since $w_k$ grows quadratically (specifically $w_k \geq |k|$ for $k\neq 0$),
only finitely many $k$ satisfy $w_k < n$.
\end{proof}

\begin{definition}
\label{def.pent.coeff}
\lean{AlgebraicCombinatorics.pentagonalCoeff}
\leanhelper
\leanok
The \emph{pentagonal coefficient} at $n\in\mathbb{N}$ is
\[
\mathrm{pentagonalCoeff}(n) =
\begin{cases}
(-1)^{|k|} & \text{if } n = w_k \text{ for some } k\in\mathbb{Z}, \\
0 & \text{otherwise}.
\end{cases}
\]
This is well-defined by the injectivity of pentagonal numbers
(Lemma~\ref{lem.pent.injective}).
\end{definition}

\begin{lemma}
\label{lem.pent.coeff-at-pentagonal}
\lean{AlgebraicCombinatorics.pentagonalCoeff_of_pentagonalNumber}
\leanhelper
\leanok
For any $k\in\mathbb{Z}$,
$\mathrm{pentagonalCoeff}(w_k) = (-1)^{|k|}$.
\end{lemma}

\begin{proof}
\leanok
Immediate from the definitions, since the inverse lookup recovers $k$ from $w_k$
by injectivity (Lemma~\ref{lem.pent.injective}).
\end{proof}

\begin{lemma}
\label{lem.pent.coeff-zero-non-pentagonal}
\lean{AlgebraicCombinatorics.pentagonalCoeff_eq_zero_of_not_pentagonal}
\leanhelper
\leanok
If $m\in\mathbb{N}$ is not a pentagonal number (i.e., $w_k \neq m$ for all $k\in\mathbb{Z}$),
then $\mathrm{pentagonalCoeff}(m) = 0$.
\end{lemma}

\begin{proof}
\leanok
By definition of pentagonal coefficients.
\end{proof}

\subsubsection{Generating function definitions}

\begin{definition}
\label{def.pent.pentagonalSeries}
\lean{AlgebraicCombinatorics.pentagonalSeries}
\leanhelper
\leanok
The \emph{pentagonal series} is the formal power series
\[
Q = \sum_{n\geq 0} \mathrm{pentagonalCoeff}(n)\, x^n
  = \sum_{k\in\mathbb{Z}} (-1)^k x^{w_k}
  \in R[[x]].
\]
\end{definition}

\begin{definition}
\label{def.pent.eulerProduct}
\lean{AlgebraicCombinatorics.eulerProduct}
\leanhelper
\leanok
The \emph{Euler product} is the infinite product
\[
\prod_{k=1}^{\infty}(1 - x^k) \in R[[x]],
\]
defined as a multipliable product in the Pi topology on $R[[x]]$.
\end{definition}

\begin{definition}
\label{def.fps.partitionGenFun}
\lean{AlgebraicCombinatorics.partitionGenFun}
\leanhelper
\leanok
The \emph{partition generating function} is
\[
P = \sum_{n\geq 0} p(n)\, x^n \in \mathbb{Z}[[x]],
\]
where $p(n)$ denotes the number of partitions of~$n$.
\end{definition}

\begin{lemma}
\label{lem.fps.partition-gen-fun-eq}
\lean{AlgebraicCombinatorics.partition_generating_function}
\leanhelper
\leanok
The partition generating function satisfies
\[
P = \sum_{n\geq 0} p(n)\, x^n = \prod_{k=1}^{\infty}\frac{1}{1-x^k}.
\]
\end{lemma}

\begin{proof}
\leanok
By Mathlib's partition generating function infrastructure.
\end{proof}

\begin{lemma}
\label{lem.fps.eulerProduct-mul-eulerProductInv}
\lean{AlgebraicCombinatorics.eulerProduct_mul_eulerProductInv}
\leanhelper
\leanok
The Euler product times the partition generating function equals $1$:
\[
\prod_{k=1}^{\infty}(1-x^k) \cdot \prod_{k=1}^{\infty}\frac{1}{1-x^k} = 1.
\]
\end{lemma}

\begin{proof}
\leanok
Each factor $(1-x^k)$ cancels with its inverse, using the per-term identity
$(1 + x^k + x^{2k} + \cdots)(1 - x^k) = 1$ and multipliability of the products.
\end{proof}

\begin{corollary}
\label{cor.pars.pn-rec}
\lean{AlgebraicCombinatorics.partition_recursive}
\leantarget
\leanok
For each positive integer $n$, we have
\begin{align*}
p\left(  n\right)   &  =\sum_{\substack{k\in\mathbb{Z};\\k\neq0}}\left(
-1\right)  ^{k-1}p\left(  n-w_{k}\right) \\
&  =p\left(  n-1\right)  +p\left(  n-2\right)  -p\left(  n-5\right)  -p\left(
n-7\right) \\
&  \ \ \ \ \ \ \ \ \ \ +p\left(  n-12\right)  +p\left(  n-15\right)  -p\left(
n-22\right)  -p\left(  n-26\right)  \pm\cdots.
\end{align*}
\end{corollary}

\begin{proof}
We have
\[
\sum_{m\in\mathbb{N}}p\left(  m\right)  x^{m}=\prod\limits_{k=1}^{\infty}\dfrac{1}{1-x^{k}}
\]
(by Theorem \ref{thm.pars.main-gf}) and
\[
\sum_{k\in\mathbb{Z}}\left(  -1\right)  ^{k}x^{w_{k}}=\prod\limits_{k=1}
^{\infty}\left(  1-x^{k}\right)
\]
(by Theorem \ref{thm.pars.pent}). Multiplying these two equalities, we obtain
\begin{align}
\left(  \sum_{m\in\mathbb{N}}p\left(  m\right)  x^{m}\right)  \cdot\left(
\sum_{k\in\mathbb{Z}}\left(  -1\right)  ^{k}x^{w_{k}}\right)   &  =\left(
\prod\limits_{k=1}^{\infty}\dfrac{1}{1-x^{k}}\right)  \cdot\left(
\prod\limits_{k=1}^{\infty}\left(  1-x^{k}\right)  \right) =1. \label{pf.cor.pars.pn-rec.1}
\end{align}

Now, let us fix a positive integer $n$. We shall compare the $x^{n}$-coefficients on both sides of (\ref{pf.cor.pars.pn-rec.1}).

The $x^{n}$-coefficient on the left hand side of (\ref{pf.cor.pars.pn-rec.1}) is
\begin{align*}
\sum_{k\in\mathbb{Z}}\left(  -1\right)  ^{k}p\left(  n-w_{k}\right)
&  =p\left(  n\right)  +\sum_{\substack{k\in\mathbb{Z};\\k\neq0}}\left(
-1\right)  ^{k}p\left(  n-w_{k}\right)  .
\end{align*}
But the $x^{n}$-coefficient on the right hand side of (\ref{pf.cor.pars.pn-rec.1})
is $0$ (since $n$ is positive). Hence, comparing the coefficients yields
\[
p\left(  n\right)  +\sum_{\substack{k\in\mathbb{Z};\\k\neq0}}\left(
-1\right)  ^{k}p\left(  n-w_{k}\right)  =0.
\]
Solving this for $p\left(  n\right)  $, we find
\[
p\left(  n\right)  =\sum_{\substack{k\in\mathbb{Z}
;\\k\neq0}}\left(  -1\right)  ^{k-1}p\left(  n-w_{k}\right)  .
\]
Corollary \ref{cor.pars.pn-rec} is thus proved.
\end{proof}

\subsection{Jacobi's triple product identity}

\subsubsection{The identity}

\begin{theorem}
[Jacobi's triple product identity, take 1]
\label{thm.pars.jtp1}
\lean{AlgebraicCombinatorics.jacobi_triple_product_fps'}
\leantarget
\leanok
In the ring
$\left(  \mathbb{Z}\left[  z^{\pm}\right]  \right)  \left[  \left[  q\right]
\right]  $, we have
\[
\prod_{n>0}\left(  \left(  1+q^{2n-1}z\right)  \left(  1+q^{2n-1}
z^{-1}\right)  \left(  1-q^{2n}\right)  \right)  =\sum_{\ell\in\mathbb{Z}
}q^{\ell^{2}}z^{\ell}.
\]
\end{theorem}

\begin{proof}
Both sides are equal after multiplying by the partition generating
function $\prod_{k>0}(1-q^{2k})^{-1}$. This is shown in
Lemma~\ref{lem.jtp.stateGenFun-eq-rhs} (for the RHS) and
Lemma~\ref{lem.jtp.stateGenFun-eq-lhs} (for the LHS). Since the
partition generating function is a unit
(Lemma~\ref{lem.jtp.partitionGenFunJacobi-isUnit}), the result follows
by cancellation.
\end{proof}

\begin{theorem}
[Jacobi's triple product identity, take 2]
\label{thm.pars.jtp2}
\lean{AlgebraicCombinatorics.jacobi_triple_product}
\leantarget
\leanok
Let $a$ and $b$
be two integers such that $a>0$ and $a\geq\left\vert b\right\vert $. Let
$u,v\in\mathbb{Q}$ be rational numbers with $v\neq0$. In the ring
$\mathbb{Q}\left[\left[  x\right]\right]  $, set $q=ux^{a}$ and $z=vx^{b}$.
Then,
\[
\prod_{n>0}\left(  \left(  1+q^{2n-1}z\right)  \left(  1+q^{2n-1}
z^{-1}\right)  \left(  1-q^{2n}\right)  \right)  =\sum_{\ell\in\mathbb{Z}
}q^{\ell^{2}}z^{\ell}.
\]
\end{theorem}

\begin{proof}
Both sides are equal after multiplying by the partition generating
function $\prod_{k>0}(1-q^{2k})^{-1}$. This is shown in
Lemma~\ref{lem.jtp.jacobiRHS-mul-partGenFun} (for the RHS) and
Lemma~\ref{lem.jtp.jacobiLHS-mul-partGenFun} (for the LHS). Since the
partition generating function is a unit
(Lemma~\ref{lem.jtp.partitionGenFunEval-isUnit}), the result follows
by cancellation.
\end{proof}

\subsubsection{Borcherds' state infrastructure}

\begin{definition}
\label{def.jtp.state}
\lean{AlgebraicCombinatorics.State}
\leanhelper
\leanok
A \emph{level} is a number of the form $p+\frac{1}{2}$ with
$p\in\mathbb{Z}$. A \emph{state} is a set of levels that contains
all but finitely many negative levels, and only finitely many positive levels.

For any state $S$, we define:
\begin{itemize}
\item the \emph{energy} of $S$ to be
\[
\operatorname{energy}S:=\sum_{\substack{p>0;\\p\in S}}2p
-\sum_{\substack{p<0;\\p\notin S}}2p \in\mathbb{N};
\]

\item the \emph{particle number} of $S$ to be
\[
\operatorname{parnum}S:=\left(\text{\# of positive levels in } S\right)
-\left(\text{\# of negative levels not in } S\right)\in\mathbb{Z}.
\]
\end{itemize}
\end{definition}

\begin{definition}
\label{def.jtp.groundState}
\lean{AlgebraicCombinatorics.groundState}
\leanhelper
\leanok
For $\ell\in\mathbb{Z}$, the \emph{$\ell$-ground state} is
\[
G_{\ell}:=\left\{  \text{all levels }<\ell\right\}.
\]
\end{definition}

\begin{theorem}
\label{thm.jtp.groundState-energy-parnum}
\lean{AlgebraicCombinatorics.groundState_energy, AlgebraicCombinatorics.groundState_parnum}
\leanhelper
\leanok
For any $\ell\in\mathbb{Z}$, we have
$\operatorname{energy}G_{\ell}=\ell^{2}$ and
$\operatorname{parnum}G_{\ell}=\ell$.
\end{theorem}

\begin{proof}
\leanok
If $\ell\geq0$, the energy is
$1+3+5+\cdots+(2\ell-1)=\ell^{2}$
and the particle number is $\ell-0=\ell$.
If $\ell<0$, the energy is
$-(-1)-(-3)-\cdots-(2\ell+1)=\ell^{2}$
and the particle number is $0-(-\ell)=\ell$.
\end{proof}

\begin{definition}
\label{def.jtp.jump}
\lean{AlgebraicCombinatorics.State.jump}
\leanhelper
\leanok
If $S$ is a state, $p\in S$, and $q$ is a positive integer such that
$p+q\notin S$, then the state
\[
\operatorname{jump}_{p,q}S:=\left(  S\setminus\left\{  p\right\}
\right)  \cup\left\{  p+q\right\}
\]
is said to be obtained from $S$ by letting the electron at level $p$
\emph{jump} $q$ steps to the right.
\end{definition}

\begin{theorem}
\label{thm.jtp.jump-parnum-energy}
\lean{AlgebraicCombinatorics.jump_parnum, AlgebraicCombinatorics.jump_energy}
\leanhelper
\leanok
A jump by $q$ steps keeps the particle number unchanged and raises the
energy by $2q$.
\end{theorem}

\begin{proof}
\leanok
Three cases are checked: the particle jumps from a positive to a
positive level, from a negative to a positive level, or from a negative
to a negative level.
\end{proof}

\begin{definition}
\label{def.jtp.excitedState}
\lean{AlgebraicCombinatorics.excitedState}
\leanhelper
\leanok
For any partition $\lambda=(\lambda_{1},\lambda_{2},\ldots,\lambda_{k})$
and any $\ell\in\mathbb{Z}$, the \emph{excited state} $E_{\ell,\lambda}$
is defined by starting with the $\ell$-ground state $G_{\ell}$ and then
successively letting the $k$ electrons at the highest levels jump
$\lambda_{1},\lambda_{2},\ldots,\lambda_{k}$ steps to the right,
respectively:
\[
E_{\ell,\lambda}=\left\{  \text{all levels }<\ell-k\right\}  \cup\left\{
\ell-i+\tfrac{1}{2}+\lambda_{i}\ \mid\ i\in\{1,\ldots,k\}\right\}.
\]
\end{definition}

\begin{lemma}
\label{lem.jtp.excitedState-reachable}
\lean{AlgebraicCombinatorics.excitedState_reachable}
\leanhelper
\leanok
The excited state $E_{\ell,\lambda}$ is reachable from the ground state $G_\ell$
by a sequence of jumps (in the sense of Definition~\ref{def.jtp.jump}).
\end{lemma}

\begin{proof}
\leanok
By induction on the number of parts: starting from $G_\ell$, one performs $k$
successive jumps of sizes $\lambda_1, \ldots, \lambda_k$. The intermediate
state construction ensures each jump is valid (source is occupied, target is vacant).
\end{proof}

\begin{theorem}
\label{thm.jtp.excitedState-energy-parnum}
\lean{AlgebraicCombinatorics.excitedState_energy, AlgebraicCombinatorics.excitedState_parnum}
\leanhelper
\leanok
The excited state $E_{\ell,\lambda}$ has energy $\ell^{2}+2|\lambda|$ and
particle number $\ell$.
\end{theorem}

\begin{proof}
\leanok
The state $E_{\ell,\lambda}$ is obtained from $G_{\ell}$ by $k$ jumps of
$\lambda_{1},\ldots,\lambda_{k}$ steps. Each jump by $q$ steps raises the
energy by $2q$ and keeps the particle number unchanged. Since the ground
state has energy $\ell^{2}$ and particle number $\ell$, the result follows.
\end{proof}

\subsubsection{The partition--state bijection}

\begin{lemma}
\label{lem.jtp.excitedState-injective}
\lean{AlgebraicCombinatorics.excitedState_injective}
\leanhelper
\leanok
For fixed $\ell\in\mathbb{Z}$, the map $\lambda\mapsto E_{\ell,\lambda}$ is
injective: if $E_{\ell,\mu_1} = E_{\ell,\mu_2}$, then $\mu_1 = \mu_2$.
\end{lemma}

\begin{proof}
\leanok
Two excited states are equal if and only if their level sets agree.
The partition $\lambda$ is recovered from the excited levels via
$\lambda_i = (\text{$i$-th largest excited level}) - (\ell - i + \frac{1}{2})$,
so equal level sets force equal partitions. The key step is that
$i\mapsto \lambda_i - i$ is strictly decreasing for sorted parts,
making the level encoding injective.
\end{proof}

\begin{lemma}
\label{lem.jtp.excitedState-surjective}
\lean{AlgebraicCombinatorics.excitedState_surjective}
\leanhelper
\leanok
For any $\ell\in\mathbb{Z}$ and any state $S$ with $\operatorname{parnum}(S) = \ell$,
there exists a partition $\mu$ such that $S = E_{\ell,\mu}$.
\end{lemma}

\begin{proof}
\leanok
Given $S$ with particle number $\ell$, we construct $\mu$ by listing the excited
levels of $S$ (those above the base level) in decreasing order and computing
$\mu_i = (\text{$i$-th excited level}) - (\ell - i + \frac{1}{2})$.
The proof shows these differences are positive integers forming a valid partition,
and that the resulting excited state equals~$S$.
\end{proof}

\begin{theorem}
\label{thm.jtp.partitionToState-bijective}
\lean{AlgebraicCombinatorics.partitionToState_bijective}
\leanhelper
\leanok
For each $\ell\in\mathbb{Z}$, the map
\begin{align*}
\Phi_{\ell}:\left\{  \text{partitions}\right\}   &  \rightarrow\left\{
\text{states with particle number }\ell\right\}  ,\\
\lambda &  \mapsto E_{\ell,\lambda}
\end{align*}
is a bijection.
\end{theorem}

\begin{proof}
Injectivity is Lemma~\ref{lem.jtp.excitedState-injective}.
Surjectivity is Lemma~\ref{lem.jtp.excitedState-surjective}.
\end{proof}

\begin{theorem}
\label{thm.jtp.intPartitionToState-bijective}
\lean{AlgebraicCombinatorics.intPartitionToState_bijective}
\leanhelper
\leanok
The global map
\begin{align*}
\Phi : \mathbb{Z} \times \{\text{partitions}\} &\to \{\text{states}\}, \\
(\ell, \mu) &\mapsto E_{\ell,\mu}
\end{align*}
is a bijection. The inverse sends a state $S$ to $(\operatorname{parnum}(S), \mu)$
where $\mu$ is the unique partition with $E_{\operatorname{parnum}(S), \mu} = S$.
\end{theorem}

\begin{proof}
Injectivity: if $E_{\ell_1,\mu_1} = E_{\ell_2,\mu_2}$, comparing particle
numbers gives $\ell_1 = \ell_2$, and then injectivity of $\Phi_\ell$ gives
$\mu_1 = \mu_2$. Surjectivity: any state $S$ has particle number $\ell = \operatorname{parnum}(S)$,
and by the surjectivity of $\Phi_\ell$, there exists $\mu$ with $E_{\ell,\mu} = S$.
\end{proof}

\begin{lemma}
\label{lem.jtp.energy-eq-parnum-sq-add-even}
\lean{AlgebraicCombinatorics.energy_eq_parnum_sq_add_even}
\leanhelper
\leanok
For any state $S$, there exists $n\geq 0$ such that
\[
\operatorname{energy}(S) = |\operatorname{parnum}(S)|^2 + 2n.
\]
In particular, the energy is always at least $|\operatorname{parnum}(S)|^2$,
and the excess is always even.
\end{lemma}

\begin{proof}
By surjectivity (Lemma~\ref{lem.jtp.excitedState-surjective}),
$S = E_{\ell,\mu}$ for some $\ell = \operatorname{parnum}(S)$ and partition $\mu$ of some $n$.
Then $\operatorname{energy}(S) = \ell^2 + 2n$ by Theorem~\ref{thm.jtp.excitedState-energy-parnum}.
\end{proof}

\subsubsection{Finset pair bijection with states}

\begin{definition}
\label{def.jtp.fromFinsetPair}
\lean{AlgebraicCombinatorics.State.fromFinsetPair}
\leanhelper
\leanok
Given a pair $(P, N)$ of finite subsets of $\mathbb{N}$, define the state
$\operatorname{fromFinsetPair}(P, N)$ whose levels are
\[
\{p\geq 0 : p \in P\} \cup \{p < 0 : (-p-1) \notin N\}.
\]
Here $P$ encodes the nonneg levels present, and $N$ encodes the
negative levels missing (with $n \in N$ meaning level $-(n+1)$ is missing).
\end{definition}

\begin{lemma}
\label{lem.jtp.fromFinsetPair-energy}
\lean{AlgebraicCombinatorics.State.fromFinsetPair_energy}
\leanhelper
\leanok
For any pair $(P, N)$ of finite subsets of $\mathbb{N}$,
\[
\operatorname{energy}(\operatorname{fromFinsetPair}(P,N))
= \sum_{n\in P}(2n+1) + \sum_{n\in N}(2n+1).
\]
\end{lemma}

\begin{proof}
For nonneg level $n\in P$: contribution is $2n+1$. For missing negative level $-(n+1)$
with $n\in N$: $|{-}(n{+}1)| = n+1$, so contribution is $2(n+1)-1 = 2n+1$. Both give $2n+1$.
\end{proof}

\begin{lemma}
\label{lem.jtp.fromFinsetPair-parnum}
\lean{AlgebraicCombinatorics.State.fromFinsetPair_parnum}
\leanhelper
\leanok
For any pair $(P, N)$ of finite subsets of $\mathbb{N}$,
\[
\operatorname{parnum}(\operatorname{fromFinsetPair}(P,N)) = |P| - |N|.
\]
\end{lemma}

\begin{proof}
\leanok
The number of nonneg levels present is $|P|$, and the number of
negative levels missing is $|N|$.
\end{proof}

\begin{theorem}
\label{thm.jtp.finsetPair-bijective}
\lean{AlgebraicCombinatorics.State.finsetPair_bijective}
\leanhelper
\leanok
The map $(P,N)\mapsto \operatorname{fromFinsetPair}(P,N)$ is a bijection
from pairs of finite subsets of $\mathbb{N}$ to states.
\end{theorem}

\begin{proof}
\leanok
Injectivity: $P$ and $N$ are recovered from any state $S$ by taking
$P = \{p \geq 0 : p \in S\}$ and $N = \{n \in \mathbb{N} : -(n+1) \notin S\}$.
Surjectivity: any state $S$ equals $\operatorname{fromFinsetPair}(P, N)$
where $P$ and $N$ are constructed as above.
\end{proof}

\begin{lemma}
\label{lem.jtp.finsetPair-sum-eq-partition-sum}
\lean{AlgebraicCombinatorics.State.finsetPair_sum_eq_partition_sum}
\leanhelper
\leanok
For any $d\in\mathbb{N}$ and function $f:\mathbb{Z}\to\mathbb{Z}[z^{\pm}]$,
\[
\sum_{\substack{(P,N):\\\operatorname{energy}=d}} f(|P|-|N|)
= \sum_{\substack{(\ell,\mu):\\\ell^2+2|\mu|=d}} f(\ell).
\]
Both sums are over finite sets, and equality holds because both sides can
be rewritten as $\sum_{S:\operatorname{energy}(S)=d} f(\operatorname{parnum}(S))$
using the two bijections (finset pairs and integer--partition pairs with states).
\end{lemma}

\begin{proof}
Both sides equal $\sum_{S:\operatorname{energy}(S)=d} f(\operatorname{parnum}(S))$
via the respective bijections.
The LHS uses the finset pair bijection (Theorem~\ref{thm.jtp.finsetPair-bijective}),
and the RHS uses the integer--partition bijection
(Theorem~\ref{thm.jtp.intPartitionToState-bijective}).
\end{proof}

\subsubsection{Jacobi ring definitions and state generating function}

\begin{definition}
\label{def.jtp.stateMonomial}
\lean{AlgebraicCombinatorics.stateMonomial}
\leanhelper
\leanok
The \emph{state monomial} with energy $e$ and particle number $p$ is
the element $q^e z^p \in (\mathbb{Z}[z^{\pm}])[[q]]$.
\end{definition}

\begin{definition}
\label{def.jtp.stateGenFun}
\lean{AlgebraicCombinatorics.stateGenFun}
\leanhelper
\leanok
The \emph{state generating function} is the formal sum over all
integer--partition pairs $(\ell, \mu)$:
\[
\mathcal{S} = \sum_{\ell\in\mathbb{Z}} \sum_{\mu\text{ partition}}
q^{\ell^2 + 2|\mu|} z^\ell
\in (\mathbb{Z}[z^{\pm}])[[q]].
\]
It is defined as the formal sum
$\sum'_{(\ell, \mu)} q^{\ell^2 + 2|\mu|} z^\ell$
over all integer--partition pairs.
\end{definition}

\begin{definition}
\label{def.jtp.partitionGenFunJacobi}
\lean{AlgebraicCombinatorics.partitionGenFunJacobi}
\leanhelper
\leanok
The \emph{partition generating function in JacobiRing} is
\[
\sum_{\mu\text{ partition}} q^{2|\mu|}
= \prod_{k>0}(1-q^{2k})^{-1}
\in (\mathbb{Z}[z^{\pm}])[[q]].
\]
\end{definition}

\begin{definition}
\label{def.jtp.jacobiZZProduct}
\lean{AlgebraicCombinatorics.jacobiZZProduct}
\leanhelper
\leanok
The \emph{$z$-dependent product} is
\[
\prod_{n>0}\left((1+q^{2n-1}z)(1+q^{2n-1}z^{-1})\right)
\in (\mathbb{Z}[z^{\pm}])[[q]].
\]
\end{definition}

\begin{definition}
\label{def.jtp.jacobiQProduct}
\lean{AlgebraicCombinatorics.jacobiQProduct}
\leanhelper
\leanok
The \emph{$q$-only product} is
\[
\prod_{n>0}(1-q^{2n})
\in (\mathbb{Z}[z^{\pm}])[[q]].
\]
\end{definition}

\subsubsection{LHS factorization and product--state connection}

\begin{lemma}
\label{lem.jtp.jacobiLHS-eq-ZZProduct-mul-QProduct}
\lean{AlgebraicCombinatorics.jacobiLHS'_eq_ZZProduct_mul_QProduct}
\leanhelper
\leanok
The LHS of Jacobi's identity factors as
\[
\prod_{n>0}\left((1+q^{2n-1}z)(1+q^{2n-1}z^{-1})(1-q^{2n})\right)
= \left(\prod_{n>0}(1+q^{2n-1}z)(1+q^{2n-1}z^{-1})\right)
\cdot \left(\prod_{n>0}(1-q^{2n})\right).
\]
\end{lemma}

\begin{proof}
\leanok
Each factor in the LHS product is
$(1+q^{2n-1}z)(1+q^{2n-1}z^{-1})(1-q^{2n})$,
which splits into the product of the first two factors (which form the $z$-dependent product)
and the third factor (which forms the $q$-only product). The infinite product
distributes because both sub-products are multipliable.
\end{proof}

\begin{lemma}
\label{lem.jtp.partitionGenFunJacobi-mul-QProduct}
\lean{AlgebraicCombinatorics.partitionGenFunJacobi_mul_QProduct_eq_one}
\leanhelper
\leanok
The partition generating function times the $q$-only product equals $1$:
\[
\prod_{k>0}(1-q^{2k})^{-1} \cdot \prod_{k>0}(1-q^{2k}) = 1.
\]
\end{lemma}

\begin{proof}
\leanok
Per-term cancellation: $(1 + q^{2k} + q^{4k} + \cdots)(1 - q^{2k}) = 1$ for each $k$.
The per-term identity and multipliability of both products give the result.
\end{proof}

\begin{lemma}
\label{lem.jtp.jacobiZZProduct-eq-double-tsum}
\lean{AlgebraicCombinatorics.jacobiZZProduct_eq_double_tsum}
\leanhelper
\leanok
The $z$-dependent product admits a binary expansion as a double sum:
\[
\prod_{n>0}((1+q^{2n-1}z)(1+q^{2n-1}z^{-1}))
= \sum_{P,N \text{ finite}} \left(\prod_{n\in P}(q^{2n+1}z)\right)\left(\prod_{n\in N}(q^{2n+1}z^{-1})\right).
\]
\end{lemma}

\begin{proof}
The product of $(1+a_n)(1+b_n)$ over $n$ expands by choosing a subset $P$ of indices
contributing $a_n$ and a subset $N$ contributing $b_n$. This uses the identities
$\prod(1+x_n) = \sum_{S\text{ finite}} \prod_{n\in S} x_n$ for each factor,
combined via the product rule for infinite products.
\end{proof}

\begin{lemma}
\label{lem.jtp.double-sum-term-explicit}
\lean{AlgebraicCombinatorics.double_sum_term_explicit}
\leanhelper
\leanok
Each term of the binary expansion has the explicit form
\[
\prod_{n\in P}(q^{2n+1}z) \cdot \prod_{n\in N}(q^{2n+1}z^{-1})
= q^{\sum_{n\in P}(2n+1) + \sum_{n\in N}(2n+1)} \cdot z^{|P| - |N|}.
\]
\end{lemma}

\begin{proof}
The $q$-exponents add up, and the $z$-powers contribute $+1$ for each element of $P$
and $-1$ for each element of $N$, giving $z^{|P|-|N|}$.
\end{proof}

\begin{lemma}
\label{lem.jtp.coeff-double-sum-eq-coeff-stateGenFun}
\lean{AlgebraicCombinatorics.coeff_double_sum_eq_coeff_stateGenFun}
\leanhelper
\leanok
For each $d\in\mathbb{N}$, the $q^d$-coefficient of the double sum
equals the $q^d$-coefficient of the state generating function:
\[
[q^d]\sum_{P,N}\prod_{n\in P}(q^{2n+1}z)\prod_{n\in N}(q^{2n+1}z^{-1})
= [q^d]\sum_{(\ell,\mu)} q^{\ell^2+2|\mu|}z^\ell.
\]
\end{lemma}

\begin{proof}
By Lemma~\ref{lem.jtp.double-sum-term-explicit}, the LHS coefficient
is a sum of $z^{|P|-|N|}$ over pairs with $\sum(2n+1) = d$.
By Lemma~\ref{lem.jtp.finsetPair-sum-eq-partition-sum}, this equals the
sum of $z^\ell$ over $(\ell,\mu)$ with $\ell^2 + 2|\mu| = d$.
\end{proof}

\begin{lemma}
\label{lem.jtp.jacobiZZProduct-eq-stateGenFun}
\lean{AlgebraicCombinatorics.jacobiZZProduct_eq_stateGenFun}
\leanhelper
\leanok
The $z$-dependent product equals the state generating function:
\[
\prod_{n>0}((1+q^{2n-1}z)(1+q^{2n-1}z^{-1})) = \mathcal{S}.
\]
\end{lemma}

\begin{proof}
\leanok
By Lemma~\ref{lem.jtp.jacobiZZProduct-eq-double-tsum}, the LHS equals the double sum.
By power series extensionality and Lemma~\ref{lem.jtp.coeff-double-sum-eq-coeff-stateGenFun},
coefficients match those of the state generating function.
\end{proof}

\subsubsection{Proof infrastructure for Jacobi's triple product identity}

\begin{lemma}
\label{lem.jtp.partitionGenFunJacobi-isUnit}
\lean{AlgebraicCombinatorics.partitionGenFunJacobi_isUnit}
\leanhelper
\leanok
The partition generating function
$\prod_{k>0}(1-q^{2k})^{-1}\in(\mathbb{Z}[z^{\pm}])[[q]]$
is a unit (its constant term is $1$).
\end{lemma}

\begin{proof}
\leanok
The constant coefficient is $1$, which is a unit.
\end{proof}

\begin{lemma}
\label{lem.jtp.stateGenFun-eq-rhs}
\lean{AlgebraicCombinatorics.stateGenFun_eq_jacobiRHS'_mul_partitionGenFunJacobi}
\leanhelper
\leanok
The state generating function equals the RHS of Jacobi's identity
times the partition generating function:
\[
\sum_{S\text{ state}}q^{\operatorname{energy}S}z^{\operatorname{parnum}S}
=\left(\sum_{\ell\in\mathbb{Z}}q^{\ell^{2}}z^{\ell}\right)
\prod_{n>0}(1-q^{2n})^{-1}.
\]
\end{lemma}

\begin{proof}
Using the bijection $\Phi_{\ell}$ from Theorem~\ref{thm.jtp.partitionToState-bijective},
we rewrite the sum over states as
\[
\sum_{\ell\in\mathbb{Z}}\sum_{\lambda\text{ partition}}
q^{\ell^{2}+2|\lambda|}z^{\ell}
=\left(\sum_{\ell\in\mathbb{Z}}q^{\ell^{2}}z^{\ell}\right)
\left(\sum_{\lambda\text{ partition}}q^{2|\lambda|}\right).
\]
The second factor is $\prod_{n>0}(1-q^{2n})^{-1}$ by the generating
function for partitions.
\end{proof}

\begin{lemma}
\label{lem.jtp.stateGenFun-eq-lhs}
\lean{AlgebraicCombinatorics.stateGenFun_eq_jacobiLHS'_mul_partitionGenFunJacobi}
\leanhelper
\leanok
The state generating function equals the LHS of Jacobi's identity
times the partition generating function:
\[
\sum_{S\text{ state}}q^{\operatorname{energy}S}z^{\operatorname{parnum}S}
=\prod_{n>0}\left((1+q^{2n-1}z)(1+q^{2n-1}z^{-1})(1-q^{2n})\right)
\cdot\prod_{k>0}(1-q^{2k})^{-1}.
\]
\end{lemma}

\begin{proof}
The proof chains the identities.
By Lemma~\ref{lem.jtp.jacobiZZProduct-eq-stateGenFun},
\[
\mathcal{S} = \prod_{n>0}((1+q^{2n-1}z)(1+q^{2n-1}z^{-1})).
\]
By Lemma~\ref{lem.jtp.partitionGenFunJacobi-mul-QProduct}, this equals
\[
\prod_{n>0}((1+q^{2n-1}z)(1+q^{2n-1}z^{-1})) \cdot \left(\prod_{n>0}(1-q^{2n}) \cdot \prod_{n>0}(1-q^{2n})^{-1}\right).
\]
By Lemma~\ref{lem.jtp.jacobiLHS-eq-ZZProduct-mul-QProduct}, this equals
\[
\prod_{n>0}\left((1+q^{2n-1}z)(1+q^{2n-1}z^{-1})(1-q^{2n})\right) \cdot \prod_{n>0}(1-q^{2n})^{-1}.
\]
\end{proof}

\subsubsection{Parameterized Jacobi identity infrastructure}

\begin{lemma}
\label{lem.jtp.partitionGenFunEval-isUnit}
\lean{AlgebraicCombinatorics.partitionGenFunEval_isUnit}
\leanhelper
\leanok
For integers $a>0$ and rational $u$, the evaluated partition generating
function $\prod_{k>0}(1-u^{2k}x^{2ak})^{-1}\in\mathbb{Q}[[x]]$
is a unit (its constant term is $1$).
\end{lemma}

\begin{proof}
\leanok
The constant coefficient is $1$, which is a unit in $\mathbb{Q}$.
\end{proof}

\begin{lemma}
\label{lem.jtp.partitionGenFunEval-mul-eulerProductParam}
\lean{AlgebraicCombinatorics.partitionGenFunEval_mul_eulerProductParam}
\leanhelper
\leanok
For $a > 0$, the partition generating function times the Euler product
(both evaluated at parameter $a$, $u$) equals $1$:
\[
\prod_{k>0}\frac{1}{1-u^{2k}x^{2ak}} \cdot \prod_{k>0}(1-u^{2k}x^{2ak}) = 1.
\]
\end{lemma}

\begin{proof}
Per-factor cancellation: each geometric series factor
$(1 + u^{2k}x^{2ak} + u^{4k}x^{4ak} + \cdots)$ multiplied by $(1 - u^{2k}x^{2ak})$ gives $1$.
The proof stabilizes the partial products and shows the coefficients eventually agree.
\end{proof}

\begin{lemma}
\label{lem.jtp.jacobiRHS-mul-partGenFun}
\lean{AlgebraicCombinatorics.jacobiRHS_mul_partitionGenFun}
\leanhelper
\leanok
The product of the RHS of Jacobi's identity (evaluated at $q=ux^a$, $z=vx^b$) with
the partition generating function equals the state generating function:
\[
\left(\sum_{\ell\in\mathbb{Z}}q^{\ell^{2}}z^{\ell}\right)
\prod_{n>0}(1-q^{2n})^{-1}
=\sum_{(\ell,\mu)}u^{\ell^{2}+2|\mu|}v^{\ell}\,x^{a(\ell^{2}+2|\mu|)+b\ell}.
\]
\end{lemma}

\begin{proof}
Follows from the Cauchy product formula for infinite sums.
\end{proof}

\begin{lemma}
\label{lem.jtp.jacobiLHS-mul-partGenFun}
\lean{AlgebraicCombinatorics.jacobiLHS_mul_partitionGenFun}
\leanhelper
\leanok
The product of the LHS of Jacobi's identity (evaluated at $q=ux^a$, $z=vx^b$) with
the partition generating function equals the state generating function.
\end{lemma}

\begin{proof}
\leanok
The proof requires:
(1) showing that the $(1-q^{2n})$ factors cancel with the partition
generating function (via Lemma~\ref{lem.jtp.partitionGenFunEval-mul-eulerProductParam}),
(2) applying binary expansion to the remaining product,
(3) reindexing using the state bijection
(Theorem~\ref{thm.jtp.partitionToState-bijective}).
\end{proof}

\subsubsection{Jacobi implies Euler}

\begin{lemma}
\label{lem.jtp.pentagonal-exponent-identity}
\lean{AlgebraicCombinatorics.pentagonal_exponent_identity}
\leanhelper
\leanok
For any $\ell\in\mathbb{Z}$,
\[
3\ell^{2}+\ell = 2\,w_{-\ell}.
\]
\end{lemma}

\begin{proof}
\leanok
Direct computation: $3\ell^{2}+\ell = (3\ell+1)\ell = (3(-\ell)-1)(-\ell) = 2w_{-\ell}$.
\end{proof}

\begin{lemma}
\label{lem.fps.fxx=gxx}
\lean{AlgebraicCombinatorics.fps_eq_of_sq_eq}
\leantarget
\leanok
Let $K$ be a commutative ring. Let $f$ and $g$ be two
FPSs in $K[[x]]$. Assume that $f[x^{2}]=g[x^{2}]$.
Then, $f=g$.
\end{lemma}

\begin{proof}
\leanok
Write $f=\sum_{n\in\mathbb{N}}f_{n}x^{n}$ and
$g=\sum_{n\in\mathbb{N}}g_{n}x^{n}$. Then
$f[x^{2}]=\sum_{n\in\mathbb{N}}f_{n}x^{2n}$ and
$g[x^{2}]=\sum_{n\in\mathbb{N}}g_{n}x^{2n}$. Comparing
$x^{2n}$-coefficients in $f[x^{2}]=g[x^{2}]$ yields $f_{n}=g_{n}$
for each $n\in\mathbb{N}$. Hence $f=g$.
\end{proof}

\begin{lemma}
\label{lem.jtp.jacobiLHS-eq-eulerProduct-sq}
\lean{AlgebraicCombinatorics.jacobiLHSEval_3_1_1_neg1_eq_eulerProduct_sq}
\leanhelper
\leanok
The LHS of Jacobi's identity evaluated at $a=3$, $b=1$, $u=1$, $v=-1$
equals $\prod_{k>0}(1-x^{k})[x^{2}]$, i.e.\ the Euler product with
$x$ replaced by $x^{2}$:
\[
\prod_{n>0}\left((1+q^{2n-1}z)(1+q^{2n-1}z^{-1})(1-q^{2n})\right)\Big|_{q=x^3,\,z=-x}
= \left(\prod_{k>0}(1-x^{k})\right)\!\big[x^{2}\big].
\]
\end{lemma}

\begin{proof}
The LHS factors simplify: setting $q=x^{3}$, $z=-x$ gives factors
$(1-x^{6n-2})(1-x^{6n-4})(1-x^{6n})$ for each $n>0$, which is
$\prod_{k>0}(1-(x^{2})^{k})$.
\end{proof}

\begin{lemma}
\label{lem.jtp.jacobiRHS-eq-pentagonalSeries-sq}
\lean{AlgebraicCombinatorics.jacobiRHSEval_3_1_1_neg1_eq_pentagonalSeries_sq}
\leanhelper
\leanok
The RHS of Jacobi's identity evaluated at $a=3$, $b=1$, $u=1$, $v=-1$
equals the pentagonal series with $x$ replaced by $x^{2}$:
\[
\sum_{\ell\in\mathbb{Z}}q^{\ell^2}z^{\ell}\Big|_{q=x^3,\,z=-x}
= \left(\sum_{k\in\mathbb{Z}}(-1)^k x^{w_k}\right)\!\big[x^{2}\big].
\]
\end{lemma}

\begin{proof}
The RHS sum is $\sum_{\ell\in\mathbb{Z}}(-1)^{\ell}x^{3\ell^{2}+\ell}$.
By Lemma~\ref{lem.jtp.pentagonal-exponent-identity},
$3\ell^{2}+\ell=2w_{-\ell}$, so the exponents are all even and the
sum equals $\sum_{k\in\mathbb{Z}}(-1)^{k}(x^{2})^{w_{k}}$, which is
the pentagonal series evaluated at $x^{2}$.
\end{proof}

\begin{lemma}
\label{lem.jtp.euler-pentagonal-rat}
\lean{AlgebraicCombinatorics.euler_pentagonal_number_theorem_rat}
\leanhelper
\leanok
Euler's pentagonal number theorem holds over $\mathbb{Q}$:
\[
\prod_{k=1}^{\infty}(1-x^k) = \sum_{k\in\mathbb{Z}}(-1)^k x^{w_k}
\quad\text{in } \mathbb{Q}[[x]].
\]
\end{lemma}

\begin{proof}
Apply Lemma~\ref{lem.fps.fxx=gxx} to reduce to showing $f[x^2] = g[x^2]$.
By Lemma~\ref{lem.jtp.jacobiLHS-eq-eulerProduct-sq} and
Lemma~\ref{lem.jtp.jacobiRHS-eq-pentagonalSeries-sq}, both sides squared equal the
LHS and RHS of Jacobi's identity evaluated at $a=3$, $b=1$, $u=1$, $v=-1$,
which are equal by Theorem~\ref{thm.pars.jtp2}.
\end{proof}

\begin{lemma}
\label{lem.jtp.eulerProduct-coeff-map}
\lean{AlgebraicCombinatorics.eulerProduct_coeff_map}
\leanhelper
\leanok
The coefficients of the Euler product are preserved by the ring homomorphism
$\mathbb{Z}\to\mathbb{Q}$:
\[
\iota\left([x^n]\prod_{k>0}(1-x^k)_{\mathbb{Z}}\right)
= [x^n]\prod_{k>0}(1-x^k)_{\mathbb{Q}},
\]
where $\iota:\mathbb{Z}\hookrightarrow\mathbb{Q}$ is the inclusion.
\end{lemma}

\begin{proof}
The ring homomorphism $\mathbb{Z}\to\mathbb{Q}$ sends the $\mathbb{Z}$-product to
the $\mathbb{Q}$-product by continuity of the induced map on power series, and
coefficients commute with the map.
\end{proof}

\subsubsection{Application: A recursion for the sum of divisors}

\begin{lemma}
\label{lem.jtp.partitionGenFun-mul-pentagonalSeries}
\lean{AlgebraicCombinatorics.partitionGenFun_mul_pentagonalSeries}
\leanhelper
\leanok
We have $P\cdot Q = 1$, where $P=\sum_{n\geq 0}p(n)x^{n}$ is the
partition generating function and
$Q=\sum_{k\in\mathbb{Z}}(-1)^{k}x^{w_{k}}$ is the pentagonal series.
\end{lemma}

\begin{proof}
\leanok
This follows from the partition generating function identity
$P=\prod_{k\geq 1}(1-x^{k})^{-1}$ (Lemma~\ref{lem.fps.partition-gen-fun-eq})
and Euler's pentagonal number theorem $Q=\prod_{k\geq 1}(1-x^{k})$
(Theorem~\ref{thm.pars.pent}). Their product $PQ = 1$ is established
by Lemma~\ref{lem.fps.eulerProduct-mul-eulerProductInv}.
\end{proof}

\begin{definition}
\label{def.jtp.sigmaSeries}
\lean{AlgebraicCombinatorics.sigmaSeries}
\leanhelper
\leanok
The \emph{sum-of-divisors generating function} is
\[
S = \sum_{k>0} \sigma(k)\, x^k \in \mathbb{Z}[[x]],
\]
where $\sigma(k) = \sigma_1(k) = \sum_{d\mid k} d$ is the sum of
positive divisors of~$k$.
\end{definition}

\begin{lemma}
\label{lem.jtp.partition-sigma-identity}
\lean{AlgebraicCombinatorics.partition_sigma_identity}
\leanhelper
\leanok
For each $n\in\mathbb{N}$,
\[
n\cdot p(n) = \sum_{k=1}^{n} \sigma(k)\cdot p(n-k).
\]
\end{lemma}

\begin{proof}
\leanok
The proof uses a combinatorial double-counting argument: a marked partition $(p, i)$
(a partition $p$ of $n$ together with a marked cell $i\in\{1,\ldots,n\}$)
can equivalently be described by $(d, m, j, \mu)$ where $d$ is the row length of the
marked cell, $m$ is the number of rows of length $d$, $j$ is the column of the cell,
and $\mu$ is the partition obtained by removing all rows of length $d$.
The identity follows by summing $\sum_{k=1}^n \sigma(k) p(n-k) = \sum_{d\mid k} d \cdot p(n-k)$
and reindexing.
\end{proof}

\begin{lemma}
\label{lem.jtp.X-mul-deriv-partitionGenFun}
\lean{AlgebraicCombinatorics.X_mul_deriv_partitionGenFun_eq}
\leanhelper
\leanok
We have $xP'=SP$, where $P$ is the partition generating function and
$S=\sum_{k>0}\sigma(k)x^{k}$ is the sum-of-divisors generating function.
\end{lemma}

\begin{proof}
\leanok
Both sides have the same coefficients: the coefficient of $x^{n}$ on
the left is $n\cdot p(n)$, and on the right is
$\sum_{k=1}^{n}\sigma(k)\cdot p(n-k)$. These are equal by the
combinatorial identity (Lemma~\ref{lem.jtp.partition-sigma-identity}).
\end{proof}

\begin{lemma}
\label{lem.jtp.pentagonal-deriv-identity}
\lean{AlgebraicCombinatorics.pentagonal_deriv_identity}
\leanhelper
\leanok
We have $xQ'=-QS$, where $Q$ is the pentagonal series and $S$ is the
sum-of-divisors generating function.
\end{lemma}

\begin{proof}
\leanok
From $PQ=1$, differentiate to get $P'Q+PQ'=0$, so $PQ'=-P'Q$.
Multiply by $x$: $xPQ'=-xP'Q=-SPQ=-S$ (using $xP'=SP$ and $PQ=1$).
Since $PQ=1$ is a unit, multiply both sides by $Q$ to get $xQ'=-QS$.
\end{proof}

\begin{lemma}
\label{lem.jtp.coeff-X-mul-deriv-pentagonal-match}
\lean{AlgebraicCombinatorics.coeff_X_mul_deriv_pentagonal_match}
\leanhelper
\leanok
The $n$-th coefficient of $xQ'$ is determined by pentagonal numbers:
\[
[x^n](xQ') =
\begin{cases}
(-1)^{|k|}\cdot n & \text{if } n = w_k \text{ for some } k\in\mathbb{Z}, \\
0 & \text{otherwise}.
\end{cases}
\]
This follows from $xQ' = \sum_{k\in\mathbb{Z}} (-1)^k w_k x^{w_k}$.
\end{lemma}

\begin{proof}
\leanok
Taking the derivative of $Q = \sum_n \mathrm{pentagonalCoeff}(n) x^n$ and multiplying by $x$
gives $[x^n](xQ') = n\cdot\mathrm{pentagonalCoeff}(n)$. The result follows from
Lemma~\ref{lem.pent.coeff-at-pentagonal} and Lemma~\ref{lem.pent.coeff-zero-non-pentagonal}.
\end{proof}

\begin{theorem}
\label{thm.pars.euler-sum-div-rec}
\lean{AlgebraicCombinatorics.euler_sum_divisors_recursive}
\leantarget
\leanok
For any positive integer $n$, let
$\sigma(n)$ denote the sum of all positive divisors of $n$.

Then, for each positive integer $n$, we have
\[
\sum_{\substack{k\in\mathbb{Z};\\w_{k}<n}}\left(  -1\right)  ^{k}\sigma\left(
n-w_{k}\right)  =
\begin{cases}
\left(  -1\right)  ^{k-1}n, & \text{if }n=w_{k}\text{ for some }k\in
\mathbb{Z};\\
0, & \text{if not.}
\end{cases}
\]
(Here, $w_{k}$ is the $k$-th pentagonal number, defined in Definition
\ref{def.pars.pent-num}.)
\end{theorem}

\begin{proof}
From Lemma~\ref{lem.jtp.pentagonal-deriv-identity}, we have
$xQ'=-QS$, where $Q=\sum_{k\in\mathbb{Z}}(-1)^{k}x^{w_{k}}$
and $S=\sum_{k>0}\sigma(k)x^{k}$.

Taking the derivative:
\[
xQ' = \sum_{k\in\mathbb{Z}}(-1)^{k}w_{k}x^{w_{k}}.
\]
On the other hand:
\[
-QS = -\left(\sum_{k\in\mathbb{Z}}(-1)^{k}x^{w_{k}}\right)
\left(\sum_{i>0}\sigma(i)x^{i}\right)
= \sum_{m>0}\sum_{\substack{k\in\mathbb{Z};\\m>w_{k}}}
(-1)^{k-1}\sigma(m-w_{k})x^{m}.
\]

Comparing the coefficient of $x^{n}$ on both sides:
\begin{itemize}
\item On the left, the coefficient is $(-1)^{k}w_{k}$ if $n=w_{k}$
for some $k$, and $0$ otherwise (Lemma~\ref{lem.jtp.coeff-X-mul-deriv-pentagonal-match});
\item On the right, the coefficient is
$\sum_{\substack{k\in\mathbb{Z};\\w_{k}<n}}(-1)^{k-1}\sigma(n-w_{k})$.
\end{itemize}
Rearranging gives the result.
\end{proof}

\subsubsection{Alternative proof via evaluation homomorphism}

\begin{definition}
\label{def.jtp.evalJacobiCorrect}
\lean{AlgebraicCombinatorics.evalJacobiCorrect}
\leanhelper
\leanok
The \emph{evaluation map} from $(\mathbb{Z}[z^{\pm}])[[q]]$ to $\mathbb{Q}[[x]]$
is defined by sending a power series $f = \sum_e c_e q^e$ (where $c_e \in \mathbb{Z}[z^{\pm}]$)
to
\[
\operatorname{eval}_{a,b,u,v}(f) = \sum_{e\geq 0} u^e \cdot \widehat{c_e}(a,b,v,e),
\]
where $\widehat{c_e}(a,b,v,e)$ denotes the shifted evaluation that replaces each
$z^\ell$ in $c_e$ by $v^\ell x^{ae+b\ell}$.
\end{definition}

\begin{lemma}
\label{lem.jtp.evalJacobiCorrect-jacobiRHS}
\lean{AlgebraicCombinatorics.evalJacobiCorrect_jacobiRHS'}
\leanhelper
\leanok
Evaluating the RHS of Jacobi's identity gives the parameterized RHS:
\[
\operatorname{eval}_{a,b,u,v}\!\left(\sum_{\ell\in\mathbb{Z}}q^{\ell^2}z^\ell\right)
= \sum_{\ell\in\mathbb{Z}} u^{\ell^2} v^\ell\, x^{a\ell^2+b\ell}.
\]
\end{lemma}

\begin{proof}
\leanok
Each sum term $q^{\ell^2}z^\ell$ is evaluated to $u^{\ell^2} v^\ell x^{a\ell^2+b\ell}$.
The evaluation distributes over the infinite sum by summability and continuity.
\end{proof}

\begin{theorem}
\label{thm.jtp.jacobi-via-evalCorrect}
\lean{AlgebraicCombinatorics.jacobi_triple_product_via_evalCorrect}
\leanhelper
\leanok
The parameterized Jacobi triple product identity (Theorem~\ref{thm.pars.jtp2})
can also be derived from Theorem~\ref{thm.pars.jtp1} by applying the
evaluation map $\operatorname{eval}_{a,b,u,v}$.
\end{theorem}

\begin{proof}
\leanok
Apply $\operatorname{eval}_{a,b,u,v}$ to both sides of
the identity from Theorem~\ref{thm.pars.jtp1} (i.e., the LHS $=$ RHS in
$(\mathbb{Z}[z^{\pm}])[[q]]$).
By Lemma~\ref{lem.jtp.evalJacobiCorrect-jacobiRHS}, the RHS evaluates to
the parameterized RHS. The LHS evaluates to the parameterized LHS
by a corresponding evaluation lemma.
\end{proof}
